\section{Target DFA energy evaluated on the HF density}
\label{sec:fixed-density-energy}

\subsection{The source HF problem}

The source orbitals are obtained by solving the periodic RHF/KRHF equation
self-consistently,
\begin{equation}
  F_{\mathrm{HF}}^{\vk}C^{\vk}
  =S^{\vk}C^{\vk}\varepsilon^{\vk}.
  \label{eq:hf-roothaan}
\end{equation}
The restricted HF Fock matrix is
\begin{equation}
  F_{\mathrm{HF}}^{\vk}
  =h^{\vk}+J^{\vk}[D]-\frac12K^{\vk}[D],
  \label{eq:hf-fock}
\end{equation}
where $D$ is the density of Eq.~\eqref{eq:restricted-density}, including the
occupation factor of two.  The HF energy per primitive cell is
\begin{equation}
  E_{\mathrm{HF}}
  =\frac{1}{N_k}\sum_{\vk}\Tr[D^{\vk}h^{\vk}]
   +\frac{1}{2N_k}\sum_{\vk}\Tr\!\left[
      D^{\vk}\left(J^{\vk}-\frac12K^{\vk}\right)\right]
   +E_{\mathrm{NN}}.
  \label{eq:hf-energy}
\end{equation}
FFTDF, GDF, or RSDF determines the numerical representation of $J$ and $K$ in
the source HF calculation.

\subsection{Target energy at a frozen HF density}

After the source HF calculation has converged, a target RKS/KRKS object is
constructed and the target functional is evaluated while the density remains
fixed.  For a semilocal target DFA,
\begin{equation}
  E_{\mathrm{DC-DFA}}[D]
  =E_1[D]+E_{\mathrm H}[D]+E_{\mathrm{xc}}^{\mathrm{DFA}}[D]
   +E_{\mathrm{NN}},
  \label{eq:dc-energy-components}
\end{equation}
with
\begin{align}
  E_1[D]
  &=\frac{1}{N_k}\sum_{\vk}\Tr[D^{\vk}h^{\vk}],
  \label{eq:dc-one-electron-energy}\\
  E_{\mathrm H}[D]
  &=\frac{1}{2N_k}\sum_{\vk}\Tr[D^{\vk}J^{\vk}[D]].
  \label{eq:dc-hartree-energy}
\end{align}
PBE, SCAN, and r2SCAN are pure semilocal target functionals, so the target
energy contains no exact-exchange matrix $K$.  Nevertheless, because $D$ is an
HF solution, source HF exchange remains in the orbital Hessian that determines
the response.

The target Fock matrix, namely the density derivative of
$E_{\mathrm{DC-DFA}}$, is
\begin{equation}
  F_{\mathrm t}^{\vk}
  =h^{\vk}+J^{\vk}[D]+V_{\mathrm{xc}}^{\vk}[D].
  \label{eq:target-fock}
\end{equation}
The subscript $\mathrm t$ denotes ``target.''  Once $D_{\mathrm{HF}}$ has been
passed to the target object, no KS orbital optimization is performed.

\subsection{Density and semilocal variables}

The real-space density is
\begin{equation}
  \rho(\mathbf r)
  =\frac{1}{N_k}\sum_{\vk}\sum_{\mu\nu}
    D_{\mu\nu}^{\vk}
    \phi_{\nu\vk}^{*}(\mathbf r)
    \phi_{\mu\vk}(\mathbf r).
  \label{eq:real-space-density}
\end{equation}
A spin-unpolarized GGA may be written as
\begin{equation}
  E_{\mathrm{xc}}^{\mathrm{GGA}}
  =\int_\Omega
    f\!\left(\rho,\sigma\right)\dd\mathbf r,
  \qquad
  \sigma=|\nabla\rho|^2.
  \label{eq:gga-energy}
\end{equation}

SCAN and r2SCAN also depend on the positive kinetic-energy density
\begin{equation}
  \tau(\mathbf r)
  =\frac{1}{2N_k}\sum_{\vk}\sum_{\mu\nu}
   D_{\mu\nu}^{\vk}
   \nabla\phi_{\nu\vk}^{*}(\mathbf r)\cdot
   \nabla\phi_{\mu\vk}(\mathbf r),
  \label{eq:tau-definition}
\end{equation}
so that
\begin{equation}
  E_{\mathrm{xc}}^{\mathrm{MGGA}}
  =\int_\Omega f(\rho,\sigma,\tau)\dd\mathbf r.
  \label{eq:mgga-energy}
\end{equation}
The PySCF SCAN/r2SCAN route does not use a density-Laplacian channel.  The
factor $1/2$ in Eq.~\eqref{eq:tau-definition} is part of the convention and
must be included exactly once in the contraction with
$v_\tau=\partial f/\partial\tau$.

\subsection{Nonvariational orbital gradient}

Apply an occupied--virtual anti-Hermitian rotation to the MO coefficients,
\begin{equation}
  C^{\vk}\longrightarrow C^{\vk}e^{\kappa^{\vk}},
  \qquad \kappa^{\vk\dagger}=-\kappa^{\vk}.
  \label{eq:orbital-rotation}
\end{equation}
Including the conventional closed-shell factor of two, the first variation of
the energy may be written as
\begin{equation}
  \delta E_{\mathrm{DC-DFA}}
  =\frac{2}{N_k}\operatorname{Re}
   \sum_{\vk ai}F_{\mathrm t,ai}^{\vk}\,
   \delta\kappa_{ai}^{\vk}.
  \label{eq:target-orbital-gradient}
\end{equation}
Canonical source HF orbitals satisfy $F_{\mathrm{HF},ai}^{\vk}=0$, but in
general
\begin{equation}
  F_{\mathrm t,ai}^{\vk}\ne0.
  \label{eq:target-vo-nonzero}
\end{equation}
This nonzero target virtual--occupied block is the right-hand side of the
Z-vector equation.

\subsection{Metadata required by the energy object}

Even for the same named DFA, different bases, pseudopotentials, cells,
k-points, DF backends, and numerical grids define different finite energies.
A reproducible energy object must record at least the target functional, source
reference, $N_k$, lattice and coordinate units, primary and auxiliary bases,
GTH pseudopotential, DF backend, and either the FFT mesh or the RSDF source and
metric parameters.  The gradient must differentiate precisely this same
energy object and its resolved settings.
